La creación de un examen sigue un proceso de configuración progresiva
en el que el coordinador define, por capas, todos los elementos
necesarios para que el sistema lo procese automáticamente:

\begin{enumerate}

  \item \textbf{Examen y modelos.} Se crea una entidad \texttt{Exam}
  con datos académicos básicos (asignatura, fecha, tipo) y, sobre ella,
  uno o varios \texttt{ExamModel}. Cada modelo es una variante
  estructural del mismo examen (típicamente, modelos A y B con el
  mismo contenido pero distinta disposición o problemas equivalentes
  barajados).

  \item \textbf{Estructura de páginas.} Para cada modelo se definen
  \texttt{PageProfile} que describen el contenido reconocible de cada
  página: \texttt{RecognitionZone} que enumeran las regiones donde
  aparecerán códigos \acs{qr}, casillas de marca, identificación del
  alumno o texto \acs{ocr}.

  \item \textbf{Problemas y rúbricas.} Cada examen se descompone en
  \texttt{Problem} con su puntuación máxima. Sobre cada problema se
  definen \texttt{RubricCriterion} con su descripción, peso y modo de
  aplicación.

  \item \textbf{\acs{pdf} instrumentado.} Una vez configurada la
  estructura, el sistema genera un \acs{pdf} instrumentado a partir
  del \acs{pdf} en blanco del examen: cada página recibe un código
  \acs{qr} único que codifica la organización, el examen, el modelo y
  el número de página, lo que permite el reconocimiento automático
  posterior (\cref{SEC:DES-INGEST}).

\end{enumerate}
